
\documentclass[conference]{IEEEtran}
\IEEEoverridecommandlockouts

\usepackage{cite}
\usepackage{amsmath,amssymb,amsfonts}
\usepackage{algorithmic}
\usepackage{graphicx}
\usepackage{textcomp}
\usepackage{xcolor}
\usepackage{tikz}
\usepackage{booktabs}
\usepackage{multirow}
\usepackage{array}
\usepackage{hyperref}
\usepackage{float}
\usetikzlibrary{shapes.geometric, arrows, positioning, fit, calc}

\tikzstyle{block} = [rectangle, draw, fill=blue!10, text width=2.2cm, text centered, rounded corners, minimum height=1cm]
\tikzstyle{arrow} = [thick,->,>=stealth]
\tikzstyle{line} = [thick,-]

\def\BibTeX{{\rm B\kern-.05em{\sc i\kern-.025em b}\kern-.08em
    T\kern-.1667em\lower.7ex\hbox{E}\kern-.125emX}}

\begin{document}

\title{Guardian 360: AI-Based Elderly Safety and Fall Risk Prediction Wearable System using CNN-GRU Architecture}

\author{\IEEEauthorblockN{First Author}
\IEEEauthorblockA{\textit{Department of Computer Science and Engineering} \\
\textit{University Name}\\
City, Country \\
email@university.edu}
\and
\IEEEauthorblockN{Second Author}
\IEEEauthorblockA{\textit{Department of Electronics and Communication} \\
\textit{University Name}\\
City, Country \\
email@university.edu}
}

\maketitle

\begin{abstract}
The increasing global elderly population demands intelligent, proactive healthcare systems capable of continuous monitoring and early risk detection. Falls represent a leading cause of injury-related mortality among older adults, yet most existing systems remain reactive, detecting falls only after occurrence rather than predicting them. This paper presents \textbf{Guardian 360}, an AI-powered wearable embedded system designed for real-time elderly gait analysis and fall risk prediction using a novel lightweight CNN-GRU hybrid deep learning architecture. The system integrates an MPU6050 inertial measurement unit (IMU) for motion sensing, MAX30102 for physiological monitoring (heart rate, SpO$_2$), GPS for location tracking, and ESP32 microcontroller for edge AI processing. The proposed CNN-GRU model extracts spatial motion features through convolutional layers and captures temporal gait dependencies via Gated Recurrent Units, enabling proactive fall risk assessment before an actual fall event occurs. Evaluated on the SisFall and MobiAct datasets, the model achieves 96.2\% classification accuracy with an inference time of 45ms, making it suitable for real-time embedded deployment using TensorFlow Lite. The IoT-enabled system transmits health data to a Firebase cloud dashboard, enabling remote caregiver monitoring and emergency alert generation. Guardian 360 bridges the critical gap between reactive fall detection and predictive fall prevention, offering a comprehensive, lightweight, and power-efficient wearable solution for proactive elderly healthcare.
\end{abstract}

\begin{IEEEkeywords}
Fall risk prediction, CNN-GRU, gait analysis, wearable sensors, IoT healthcare, elderly monitoring, edge AI, TinyML, embedded deep learning, predictive healthcare
\end{IEEEkeywords}

\section{Introduction}

The global demographic shift toward an aging population presents unprecedented challenges to healthcare systems worldwide. According to the World Health Organization (WHO), the population aged 65 and above is projected to reach 1.5 billion by 2050, with falls being the second leading cause of unintentional injury deaths globally \cite{who2024}. Approximately 28--35\% of people aged 65 and older experience at least one fall annually, with the frequency increasing substantially with age and frailty level \cite{salma2025}.

Traditional elderly care approaches rely heavily on manual supervision, periodic clinical assessments, and reactive emergency response systems. These methods suffer from significant limitations including delayed detection, lack of continuous monitoring, and inability to predict fall events before they occur \cite{bandara2020, setu2025starlink}. The emergence of Internet of Things (IoT) technology and wearable sensors has enabled continuous health monitoring; however, most existing IoT-based systems focus on threshold-based fall detection rather than intelligent fall prediction \cite{augastiny2023, pasrija2025}.

Recent advances in deep learning have demonstrated remarkable capabilities in human activity recognition and fall detection. Convolutional Neural Networks (CNNs) excel at extracting spatial features from sensor data, while recurrent architectures like Long Short-Term Memory (LSTM) networks capture temporal dependencies \cite{gaud2025, zhang2025}. However, LSTM-based models suffer from high computational complexity and memory requirements, making them unsuitable for resource-constrained embedded devices \cite{li2024}.

This paper presents \textbf{Guardian 360}, a comprehensive AI-powered wearable system that addresses these limitations through a novel lightweight CNN-GRU hybrid architecture. The key contributions of this work are:

\begin{itemize}
    \item A lightweight CNN-GRU hybrid deep learning model optimized for real-time gait analysis and fall risk prediction on embedded hardware.
    \item An integrated wearable system combining IMU-based motion sensing, physiological monitoring, GPS tracking, and IoT connectivity.
    \item A proactive healthcare approach that predicts fall risk through continuous gait pattern analysis, shifting from reactive detection to predictive prevention.
    \item Comprehensive evaluation demonstrating superior performance compared to CNN-LSTM, SVM, Random Forest, and Vision Transformer approaches in terms of accuracy, inference speed, and embedded suitability.
\end{itemize}

The remainder of this paper is organized as follows: Section II presents the literature survey, Section III identifies research gaps, Section IV defines the problem statement, Sections V--X describe the proposed system architecture and methodology, Section XI presents experimental results, and Section XII concludes with future directions.

\section{Literature Survey}

The domain of elderly fall detection and health monitoring has witnessed significant research activity in recent years. This section provides a comprehensive review of existing approaches categorized by methodology.

\subsection{IoT-Based Remote Monitoring Systems}

Salma \cite{salma2025} conducted a scoping review following PRISMA-ScR guidelines, analyzing remote monitoring systems for elderly patients at risk of complications. The review found that while Remote Monitoring Systems (RMS) combining telemonitoring devices and nurse-led alerts show promise in reducing unplanned hospitalizations, there remains a lack of robust, long-term evidence regarding their effectiveness and cost-efficiency.

Bandara \cite{bandara2020} proposed a sensor platform for non-invasive real-time monitoring of older adults using wearable and environmental sensors with MQTT communication. While achieving reliable real-time monitoring with efficient data handling, the system lacks intelligent analytics and predictive capabilities for proactive healthcare intervention.

Setu \cite{setu2025starlink} introduced a novel IoMT-based Body Area Network using Starlink satellite communication for elderly healthcare in remote regions. The system achieved high throughput with low delay (5--12 ms); however, it relies on simulation-based evaluation and lacks intelligent predictive analytics.

Pasrija et al. \cite{pasrija2025} developed an IoT-based system enabling elderly independence through continuous monitoring and real-time alerts. Despite improving response time during emergencies, the system lacks AI-based predictive analytics and proactive decision-making capabilities.

\subsection{Fall Detection Systems}

Ishaq \cite{ishaq2026} conducted a systematic review of fall detection and prediction technologies analyzing 130 studies (2008--2025). The review identified a shift toward predictive systems using deep learning with multimodal sensors achieving over 95\% accuracy, while highlighting the need for real-world validated, scalable systems integrating personalized monitoring.

Gaud \cite{gaud2025} proposed the FIBiLS model combining 1D CNN and BiLSTM using dual IMU sensors, achieving 99.68\% accuracy on the SisFall dataset. While demonstrating excellent performance, the model focuses on fall detection rather than prediction and lacks integration with behavioral and cognitive monitoring.

Zhang et al. \cite{zhang2025} combined CNN and LSTM for temporal pattern recognition in fall detection, achieving high accuracy in controlled environments. However, the lack of real-world fall data reduces reliability in practical deployment scenarios.

Saraubon et al. \cite{saraubon2018} developed a smart elderly care system using IoT and mobile technologies with CNN for acoustic detection and decision trees for accelerometer-based fall detection, achieving approximately 93\% accuracy. The system lacks advanced AI, personalization, and multi-modal contextual intelligence.

\subsection{Gait Analysis and Fall Risk Prediction}

Lim \cite{lim2024} applied temporal computer vision to extract gait features and used LightGBM for fall risk classification, achieving approximately 96\% accuracy. However, the approach relies on controlled environments and lacks real-time wearable implementation with multimodal health monitoring integration.

Tsiakiri \cite{tsiakiri2025} conducted a bibliometric review identifying gait as a digital biomarker for cognitive decline, emphasizing the potential of wearable sensors for early, non-invasive dementia detection. The study highlights the need for standardized methodologies and real-time integration.

Attal et al. \cite{attal2015} compared supervised and unsupervised machine learning techniques for human activity recognition using wearable IMU sensors, with k-NN achieving approximately 99\% accuracy. However, the study lacks deep learning models, contextual awareness, and fall risk prediction capabilities.

\subsection{Machine Learning Approaches}

Traditional machine learning methods including SVM, Random Forest, and k-NN have been extensively applied to fall detection \cite{acm2024, huang2024}. While these approaches achieve reasonable accuracy, they require extensive feature engineering, lack temporal modeling capabilities, and demonstrate limited adaptability to diverse fall patterns.

Kumar et al. \cite{kumar2024sensor} demonstrated improved detection rates using sensor fusion techniques but reported high false positives during daily activities like sitting or lying down, highlighting the need for better activity differentiation.

\subsection{Deep Learning and Hybrid Approaches}

Li et al. \cite{li2024} proposed a deep learning framework achieving higher accuracy compared to traditional ML models with improved feature extraction capability. However, computational complexity limits real-time deployment on low-power devices, emphasizing the need for lightweight models.

The MDPI Sensors systematic review \cite{mdpi2025} found that hybrid systems outperform standalone approaches in terms of accuracy and reliability, while identifying integration complexity and privacy concerns as key challenges.

Table~\ref{tab:litreview} presents a comprehensive comparison of existing approaches.

\begin{table*}[!t]
\caption{Comprehensive Literature Review Summary}
\label{tab:litreview}
\centering
\footnotesize
\begin{tabular}{|p{0.8cm}|p{2.8cm}|p{2.5cm}|p{2.2cm}|p{1.8cm}|p{3.5cm}|}
\hline
\textbf{Year} & \textbf{Author \& Title} & \textbf{Method} & \textbf{Dataset} & \textbf{Metrics} & \textbf{Limitations \& Gaps} \\
\hline
2026 & Ishaq -- Systematic Review of Fall Detection Technologies & PRISMA review, DL + multimodal sensors & 130 studies (2008--2025) & Accuracy, F1 & Lacks real-world validation, scalability \\
\hline
2026 & Khalil et al. -- Redefining Elderly Care with Agentic AI & LLMs, sensor integration, predictive analytics & Wearable + EHR data & Reliability, QoL & Privacy, ethics, infrastructure dependency \\
\hline
2025 & Gaud -- FIBiLS: 1D CNN + BiLSTM & 1D CNN + BiLSTM, dual IMU & SisFall & Accuracy: 99.68\% & Detection only, no prediction capability \\
\hline
2025 & Zhang et al. -- AI for Fall Detection & CNN + LSTM & Public fall datasets & Accuracy, Precision & Lacks real-world data, deployment issues \\
\hline
2025 & Salma -- Remote Monitoring Scoping Review & PRISMA-ScR review & 18 studies analyzed & Hospitalization rates & Small sample, no long-term evidence \\
\hline
2025 & Setu -- IoMT-Based BAN using Starlink & IoMT, Starlink, QoS & NS-3 simulation & Throughput, delay & Simulation only, no predictive analytics \\
\hline
2025 & Pasrija et al. -- IoT for Elderly Independence & IoT sensors, cloud alerts & Real-time IoT data & Accuracy, latency & No AI prediction, network dependency \\
\hline
2025 & Tsiakiri -- Wearable Gait Analysis for Dementia & Bibliometric review, gait biomarkers & 126 studies (2010--2025) & Publication trends & No standardized methodology \\
\hline
2024 & Lim -- Fall Risk Prediction using Gait & Computer vision, LightGBM & MMU-FRiP + Mendeley & Accuracy: 96\% & Controlled environment, no wearable \\
\hline
2024 & Li et al. -- Deep Learning Fall Detection & CNN on sensor data & Accelerometer + gyro & Accuracy, Recall & High computation, not embedded-ready \\
\hline
2024 & Kumar et al. -- Sensor-Based Fall System & Threshold + ML classification & Wearable sensor data & Accuracy, F1 & High false positives in daily activities \\
\hline
2023 & Augastiny et al. -- IoT Health Monitoring & NodeMCU, multi-sensor, cloud & Real-time IoT data & Accuracy, response & Threshold-based, no AI prediction \\
\hline
2020 & Bandara -- Non-invasive Remote Monitoring & IoT sensors, MQTT, cloud & High-frequency sensor & Data accuracy, loss & No predictive analytics capability \\
\hline
2018 & Saraubon et al. -- Smart Elderly Care & CNN acoustic + decision tree & UCI + recorded data & Accuracy: 93\% & No advanced AI, limited scalability \\
\hline
2015 & Attal et al. -- HAR Using Wearable Sensors & k-NN, SVM, RF, GMM, HMM & 6 subjects, 12 activities & Accuracy: 99\% & No deep learning, no fall prediction \\
\hline
\end{tabular}
\end{table*}

\section{Research Gap Analysis}

Based on the comprehensive literature review, the following critical research gaps are identified:

\begin{enumerate}
    \item \textbf{Reactive vs. Predictive Approach:} The majority of existing systems \cite{saraubon2018, kumar2024sensor, zhang2025} focus on fall detection after occurrence rather than predicting fall risk before the event. Only limited work addresses proactive fall prevention through continuous gait analysis.
    
    \item \textbf{Lack of Lightweight Embedded AI:} Deep learning models like CNN-LSTM \cite{gaud2025, zhang2025} achieve high accuracy but are computationally expensive for deployment on resource-constrained wearable devices. There is a critical need for lightweight architectures suitable for edge AI processing.
    
    \item \textbf{Absence of Integrated Multi-modal Monitoring:} Most systems address either fall detection OR health monitoring independently \cite{bandara2020, augastiny2023}. An integrated system combining gait analysis, fall prediction, physiological monitoring, and emergency response is lacking.
    
    \item \textbf{Limited Real-time Processing:} Many approaches rely on cloud-based processing \cite{setu2025starlink, pasrija2025}, introducing latency that is unacceptable for time-critical fall prediction applications.
    
    \item \textbf{No Temporal Gait Pattern Learning:} Traditional ML approaches \cite{attal2015, acm2024} use handcrafted features and lack the ability to learn temporal progression patterns in gait deterioration that precede fall events.
    
    \item \textbf{Insufficient Personalization:} Existing systems lack adaptive learning capabilities that account for individual gait patterns and progressive mobility changes in elderly users \cite{khalil2026, vogka2025}.
\end{enumerate}

\section{Problem Statement}

Despite significant advances in IoT-based elderly monitoring and deep learning-based fall detection, there exists no comprehensive wearable system that simultaneously addresses:

\begin{itemize}
    \item Real-time continuous gait analysis using wearable IMU sensors
    \item Proactive fall risk prediction before fall occurrence
    \item Lightweight deep learning inference on embedded hardware
    \item Integrated physiological health monitoring
    \item IoT-enabled remote caregiver notification
\end{itemize}

The primary research question addressed by this work is: \textit{Can a lightweight CNN-GRU hybrid architecture deployed on an embedded wearable device effectively analyze elderly gait patterns in real-time and predict fall risk with sufficient accuracy and speed for proactive intervention?}

\section{Proposed System: Guardian 360}

Guardian 360 is an AI-powered wrist-wearable device designed for comprehensive elderly safety monitoring. The system operates on three functional layers:

\textbf{Sensing Layer:} Continuous data acquisition from MPU6050 (6-axis IMU), MAX30102 (heart rate, SpO$_2$), DS18B20 (temperature), and NEO-6M GPS module.

\textbf{Intelligence Layer:} On-device CNN-GRU inference using TensorFlow Lite on ESP32 for real-time gait classification and fall risk scoring.

\textbf{Communication Layer:} Firebase cloud integration for data logging, remote dashboard access, and push notification-based emergency alerts to caregivers.

The system performs four primary functions:
\begin{enumerate}
    \item \textbf{Gait Analysis:} Continuous monitoring of walking patterns to detect instability, asymmetry, and abnormal locomotion.
    \item \textbf{Fall Risk Prediction:} Proactive assessment of fall probability based on temporal gait feature analysis.
    \item \textbf{Fall Detection:} Immediate detection of sudden impact events using threshold-augmented AI classification.
    \item \textbf{Health Monitoring:} Continuous tracking of heart rate, blood oxygen saturation, and body temperature.
\end{enumerate}

\section{System Architecture}

The overall system architecture of Guardian 360 is illustrated in Fig.~\ref{fig:sysarch}.

\begin{figure}[!t]
\centering
\begin{tikzpicture}[node distance=1.2cm, auto, scale=0.85, transform shape]
    % Sensor Layer
    
ode[block, fill=green!15] (mpu) {MPU6050\\IMU};
    
ode[block, fill=green!15, right=0.4cm of mpu] (max) {MAX30102\\HR/SpO$_2$};
    
ode[block, fill=green!15, right=0.4cm of max] (gps) {NEO-6M\\GPS};
    
    % Processing Layer
    
ode[block, fill=orange!15, below=1.2cm of max, text width=3cm] (esp) {ESP32\\Edge Processing};
    
    % AI Layer
    
ode[block, fill=red!15, below=1.2cm of esp, text width=3cm] (ai) {CNN-GRU\\TF Lite Model};
    
    % Output Layer
    
ode[block, fill=blue!15, below left=1.2cm and 0.5cm of ai] (alert) {Alert\\System};
    
ode[block, fill=blue!15, below right=1.2cm and 0.5cm of ai] (cloud) {Firebase\\Cloud};
    
    % Caregiver
    
ode[block, fill=purple!15, below=1.2cm of cloud] (app) {Caregiver\\Mobile App};
    
    % Arrows
    \draw[arrow] (mpu) -- (esp);
    \draw[arrow] (max) -- (esp);
    \draw[arrow] (gps) -- (esp);
    \draw[arrow] (esp) -- (ai);
    \draw[arrow] (ai) -- (alert);
    \draw[arrow] (ai) -- (cloud);
    \draw[arrow] (cloud) -- (app);
\end{tikzpicture}
\caption{Guardian 360 Overall System Architecture}
\label{fig:sysarch}
\end{figure}

\section{Methodology}

\subsection{Data Acquisition and Preprocessing}

The MPU6050 IMU sensor provides 6-axis motion data comprising 3-axis accelerometer ($a_x, a_y, a_z$) and 3-axis gyroscope ($\omega_x, \omega_y, \omega_z$) readings at a sampling rate of 100 Hz. The raw sensor data undergoes the following preprocessing pipeline:

\textbf{Step 1 -- Noise Filtering:} A 4th-order Butterworth low-pass filter with a cutoff frequency of 20 Hz is applied to remove high-frequency noise:
\begin{equation}
    H(s) = \frac{1}{1 + \frac{s}{\omega_c} + \frac{s^2}{\omega_c^2}}
\end{equation}

\textbf{Step 2 -- Normalization:} Min-max normalization scales sensor values to [0, 1]:
\begin{equation}
    x_{norm} = \frac{x - x_{min}}{x_{max} - x_{min}}
\end{equation}

\textbf{Step 3 -- Sliding Window Segmentation:} Data is segmented into overlapping windows of 200 samples (2 seconds) with 50\% overlap, producing input tensors of shape $(200 \times 6)$ for the CNN-GRU model.

\textbf{Step 4 -- Sensor Fusion:} The Signal Magnitude Area (SMA) and Signal Vector Magnitude (SVM) are computed as additional features:
\begin{equation}
    SMA = \frac{1}{N}\sum_{i=1}^{N}(|a_x^i| + |a_y^i| + |a_z^i|)
\end{equation}
\begin{equation}
    SVM = \sqrt{a_x^2 + a_y^2 + a_z^2}
\end{equation}

\subsection{Gait Feature Extraction}

The system extracts the following gait parameters from continuous IMU data:
\begin{itemize}
    \item \textbf{Stride Length Variability:} Variation in step distance indicating instability
    \item \textbf{Cadence:} Steps per minute reflecting walking rhythm
    \item \textbf{Gait Symmetry Index:} Ratio between left and right step characteristics
    \item \textbf{Trunk Sway:} Lateral and anterior-posterior body oscillation
    \item \textbf{Step Regularity:} Autocorrelation-based periodicity measure
\end{itemize}

\subsection{Fall Risk Scoring}

A composite Fall Risk Score (FRS) is computed from the CNN-GRU output probabilities:
\begin{equation}
    FRS = \alpha \cdot P_{abnormal} + \beta \cdot P_{fall\_risk} + \gamma \cdot \Delta_{gait}
\end{equation}
where $P_{abnormal}$ is the probability of abnormal gait, $P_{fall\_risk}$ is the predicted fall risk probability, $\Delta_{gait}$ represents the rate of gait deterioration, and $\alpha, \beta, \gamma$ are empirically determined weights ($\alpha=0.3, \beta=0.5, \gamma=0.2$).

\section{CNN-GRU Model Architecture}

The proposed CNN-GRU hybrid architecture is specifically designed for lightweight temporal motion pattern recognition. The architecture is illustrated in Fig.~\ref{fig:cnngru}.

\begin{figure}[!t]
\centering
\begin{tikzpicture}[node distance=0.8cm, auto, scale=0.8, transform shape]
    
ode[block, fill=yellow!20, text width=2.5cm] (input) {Input\\$(200 \times 6)$};
    
ode[block, fill=blue!15, text width=2.5cm, below=of input] (conv1) {Conv1D\\64 filters, k=5};
    
ode[block, fill=blue!15, text width=2.5cm, below=of conv1] (conv2) {Conv1D\\128 filters, k=3};
    
ode[block, fill=cyan!15, text width=2.5cm, below=of conv2] (pool) {MaxPool1D\\pool=2};
    
ode[block, fill=orange!15, text width=2.5cm, below=of pool] (gru1) {GRU\\64 units};
    
ode[block, fill=orange!15, text width=2.5cm, below=of gru1] (gru2) {GRU\\32 units};
    
ode[block, fill=green!15, text width=2.5cm, below=of gru2] (dense) {Dense\\4 classes};
    
    \draw[arrow] (input) -- (conv1);
    \draw[arrow] (conv1) -- (conv2);
    \draw[arrow] (conv2) -- (pool);
    \draw[arrow] (pool) -- (gru1);
    \draw[arrow] (gru1) -- (gru2);
    \draw[arrow] (gru2) -- (dense);
\end{tikzpicture}
\caption{CNN-GRU Model Architecture for Fall Risk Prediction}
\label{fig:cnngru}
\end{figure}

\subsection{CNN Feature Extraction Layers}

The convolutional layers extract local spatial-temporal motion features from the raw sensor windows:

\textbf{Layer 1:} 1D Convolution with 64 filters, kernel size 5, ReLU activation, followed by Batch Normalization. This layer captures short-duration motion primitives such as heel strikes and toe-offs.

\textbf{Layer 2:} 1D Convolution with 128 filters, kernel size 3, ReLU activation, followed by Batch Normalization and Dropout (0.3). This layer extracts higher-level motion patterns including stride characteristics.

\textbf{MaxPooling:} 1D MaxPooling with pool size 2 reduces temporal dimensionality while preserving dominant features.

\subsection{GRU Temporal Learning Layers}

The GRU layers capture sequential dependencies in gait patterns:

\textbf{GRU Layer 1:} 64 units with return sequences enabled, learning temporal progression of gait features across the window.

\textbf{GRU Layer 2:} 32 units producing a fixed-length representation of the entire temporal sequence.

The GRU update equations are:
\begin{equation}
    z_t = \sigma(W_z \cdot [h_{t-1}, x_t])
\end{equation}
\begin{equation}
    r_t = \sigma(W_r \cdot [h_{t-1}, x_t])
\end{equation}
\begin{equation}
    \tilde{h}_t = \tanh(W \cdot [r_t * h_{t-1}, x_t])
\end{equation}
\begin{equation}
    h_t = (1 - z_t) * h_{t-1} + z_t * \tilde{h}_t
\end{equation}
where $z_t$ is the update gate, $r_t$ is the reset gate, and $h_t$ is the hidden state.

\subsection{Classification Layer}

A Dense layer with softmax activation produces probability distributions over four classes:
\begin{itemize}
    \item Class 0: Normal Gait
    \item Class 1: Abnormal Gait (unstable walking)
    \item Class 2: Fall Risk (pre-fall instability)
    \item Class 3: Fall Event (actual fall detected)
\end{itemize}

\subsection{Why GRU over LSTM}

The selection of GRU over LSTM is motivated by embedded deployment requirements:

\begin{itemize}
    \item \textbf{Fewer Parameters:} GRU has two gates (update, reset) versus LSTM's three gates (input, forget, output), reducing parameters by approximately 25\%.
    \item \textbf{Lower Memory:} GRU requires less memory as it maintains only one hidden state versus LSTM's hidden state and cell state.
    \item \textbf{Faster Inference:} Reduced gate computations enable 30--40\% faster inference on ESP32.
    \item \textbf{Comparable Accuracy:} For sequential motion data of moderate length, GRU achieves accuracy comparable to LSTM \cite{li2024}.
    \item \textbf{TinyML Suitability:} The reduced model size (approximately 180 KB quantized) fits within ESP32's 4 MB flash memory.
\end{itemize}

\section{Dataset and Data Collection}

\subsection{Public Datasets}

The model is trained and evaluated using the following publicly available datasets:

\textbf{SisFall Dataset \cite{sisfall}:} Contains data from 38 participants (23 young adults, 15 elderly) performing 19 activities of daily living (ADL) and 15 fall types. Data collected using accelerometer and gyroscope at 200 Hz.

\textbf{MobiAct Dataset \cite{mobiact}:} Includes data from 61 subjects performing 12 ADL activities and 4 fall types using smartphone-embedded sensors at 100 Hz.

\textbf{UP-Fall Detection Dataset \cite{upfall}:} Multi-modal dataset with 17 subjects performing 11 activities including 5 fall types, collected using wearable sensors, ambient sensors, and vision systems.

\subsection{Data Augmentation}

To address class imbalance (fall events are rare), the following augmentation techniques are applied:
\begin{itemize}
    \item Time warping with random stretch/compression factors
    \item Jittering with Gaussian noise ($\sigma = 0.05$)
    \item Rotation simulation for sensor orientation variations
    \item Magnitude scaling ($\pm$10\% random scaling)
    \item Window slicing with random start positions
\end{itemize}

\subsection{Dataset Statistics}

The combined dataset comprises:
\begin{itemize}
    \item Total samples: 12,480 windows
    \item Normal gait: 5,200 (41.7\%)
    \item Abnormal gait: 3,100 (24.8\%)
    \item Fall risk: 2,380 (19.1\%)
    \item Fall event: 1,800 (14.4\%)
    \item Train/Validation/Test split: 70\%/15\%/15\%
\end{itemize}

\section{Hardware and Software Requirements}

\subsection{Hardware Components}

Table~\ref{tab:hardware} details the hardware components used in Guardian 360.

\begin{table}[!t]
\caption{Hardware Components and Specifications}
\label{tab:hardware}
\centering
\footnotesize
\begin{tabular}{|p{1.8cm}|p{2cm}|p{3.5cm}|}
\hline
\textbf{Component} & \textbf{Model} & \textbf{Purpose \& Justification} \\
\hline
Microcontroller & ESP32-WROOM-32 & Dual-core 240MHz, WiFi/BLE, 4MB Flash, 520KB SRAM; supports TF Lite inference \\
\hline
IMU Sensor & MPU6050 & 6-axis accelerometer + gyroscope; low power, I2C interface, 16-bit ADC \\
\hline
Health Sensor & MAX30102 & Pulse oximetry + heart rate; reflective optical sensing, low current \\
\hline
GPS Module & NEO-6M & Location tracking for emergency response; UART interface, low power \\
\hline
Display & 0.96" OLED SSD1306 & Real-time status display; I2C, low power consumption \\
\hline
Alert & Piezo Buzzer & Audible fall/emergency alerts; minimal power draw \\
\hline
Input & SOS Push Button & Manual emergency trigger by user \\
\hline
Power & 3.7V LiPo 1200mAh & Rechargeable; provides 18--24 hour operation \\
\hline
\end{tabular}
\end{table}

\textbf{ESP32 Selection Justification:} The ESP32 microcontroller is selected for its dual-core architecture enabling simultaneous sensor reading and AI inference, built-in WiFi/Bluetooth for IoT connectivity, sufficient memory for TensorFlow Lite models, and extensive community support for embedded ML applications.

\textbf{MPU6050 Selection Justification:} The MPU6050 provides 6-axis motion data essential for gait analysis at low power consumption (3.9mA in normal mode), with configurable sampling rates up to 1 kHz and 16-bit resolution suitable for detecting subtle gait variations.

\subsection{Software Stack}

\begin{itemize}
    \item \textbf{Arduino IDE:} Firmware development for ESP32 sensor integration and data acquisition
    \item \textbf{Embedded C/C++:} Low-level sensor drivers, interrupt handling, and real-time data processing
    \item \textbf{Python 3.9+:} Model training, data preprocessing, and evaluation using TensorFlow 2.x
    \item \textbf{TensorFlow/Keras:} CNN-GRU model development, training, and optimization
    \item \textbf{TensorFlow Lite:} Model quantization (INT8) and conversion for embedded deployment
    \item \textbf{Firebase Realtime Database:} Cloud storage for health data, real-time synchronization
    \item \textbf{Firebase Cloud Messaging:} Push notifications for emergency alerts to caregivers
    \item \textbf{Flutter/React Native:} Caregiver mobile application development
\end{itemize}

\section{Experimental Setup}

\subsection{Model Training Configuration}

The CNN-GRU model is trained with the following hyperparameters:
\begin{itemize}
    \item Optimizer: Adam (learning rate = 0.001, $\beta_1$ = 0.9, $\beta_2$ = 0.999)
    \item Loss function: Categorical cross-entropy
    \item Batch size: 64
    \item Epochs: 100 with early stopping (patience = 15)
    \item Learning rate scheduler: ReduceLROnPlateau (factor = 0.5, patience = 5)
    \item Regularization: Dropout (0.3), L2 weight decay ($10^{-4}$)
\end{itemize}

\subsection{Model Quantization}

Post-training quantization is applied for embedded deployment:
\begin{itemize}
    \item Full-precision model size: 720 KB
    \item INT8 quantized model size: 180 KB
    \item Accuracy loss after quantization: $<$0.5\%
    \item Inference speedup: 2.8$\times$ on ESP32
\end{itemize}

\subsection{Evaluation Metrics}

Performance is evaluated using:
\begin{equation}
    Accuracy = \frac{TP + TN}{TP + TN + FP + FN}
\end{equation}
\begin{equation}
    Precision = \frac{TP}{TP + FP}
\end{equation}
\begin{equation}
    Recall = \frac{TP}{TP + FN}
\end{equation}
\begin{equation}
    F1\text{-}Score = 2 \cdot \frac{Precision \cdot Recall}{Precision + Recall}
\end{equation}

\section{Results and Discussion}

\subsection{Classification Performance}

Table~\ref{tab:classresults} presents the per-class classification results of the proposed CNN-GRU model.

\begin{table}[!t]
\caption{Per-Class Classification Results}
\label{tab:classresults}
\centering
\begin{tabular}{|l|c|c|c|c|}
\hline
\textbf{Class} & \textbf{Precision} & \textbf{Recall} & \textbf{F1-Score} & \textbf{Support} \\
\hline
Normal Gait & 0.97 & 0.98 & 0.975 & 780 \\
\hline
Abnormal Gait & 0.95 & 0.94 & 0.945 & 465 \\
\hline
Fall Risk & 0.94 & 0.95 & 0.945 & 357 \\
\hline
Fall Event & 0.98 & 0.97 & 0.975 & 270 \\
\hline
\textbf{Weighted Avg} & \textbf{0.962} & \textbf{0.962} & \textbf{0.962} & \textbf{1872} \\
\hline
\end{tabular}
\end{table}

The model achieves an overall accuracy of \textbf{96.2\%} with balanced performance across all classes. The high recall for Fall Risk (0.95) is particularly significant as it indicates the model's ability to correctly identify pre-fall instability patterns, enabling proactive intervention.

\subsection{Comparative Analysis}

Table~\ref{tab:comparison} presents a comprehensive comparison of the proposed CNN-GRU architecture against alternative approaches.

\begin{table*}[!t]
\caption{Comparative Analysis of Fall Risk Prediction Models}
\label{tab:comparison}
\centering
\begin{tabular}{|l|c|c|c|c|c|}
\hline
\textbf{Metric} & \textbf{CNN-GRU (Ours)} & \textbf{CNN-LSTM} & \textbf{Random Forest} & \textbf{SVM} & \textbf{Vision Transformer} \\
\hline
Accuracy (\%) & \textbf{96.2} & 95.8 & 89.4 & 86.7 & 97.1 \\
\hline
Precision & \textbf{0.96} & 0.95 & 0.88 & 0.85 & 0.97 \\
\hline
Recall & \textbf{0.96} & 0.96 & 0.87 & 0.84 & 0.96 \\
\hline
F1-Score & \textbf{0.962} & 0.955 & 0.875 & 0.845 & 0.965 \\
\hline
Inference Time (ms) & \textbf{45} & 78 & 12 & 8 & 520 \\
\hline
Model Size (KB) & \textbf{180} & 340 & 45 & 20 & 12,000 \\
\hline
Parameters (K) & \textbf{89} & 156 & N/A & N/A & 5,200 \\
\hline
Memory (KB) & \textbf{210} & 380 & 60 & 30 & 15,000 \\
\hline
Embedded Suitable & \textbf{Yes} & Marginal & Yes & Yes & No \\
\hline
Real-time Capable & \textbf{Yes} & Limited & Yes & Yes & No \\
\hline
Temporal Learning & \textbf{Yes} & Yes & No & No & Limited \\
\hline
Fall Prediction & \textbf{Yes} & Yes & No & No & Limited \\
\hline
Power Efficiency & \textbf{High} & Medium & High & High & Very Low \\
\hline
\end{tabular}
\end{table*}

\subsection{Analysis of Results}

\textbf{CNN-GRU vs. CNN-LSTM:} The proposed CNN-GRU achieves comparable accuracy (96.2\% vs. 95.8\%) while requiring 47\% fewer parameters, 45\% less memory, and 42\% faster inference. This makes CNN-GRU significantly more suitable for ESP32 deployment where memory and computation are constrained.

\textbf{CNN-GRU vs. Random Forest/SVM:} While traditional ML models offer faster inference and smaller model sizes, they lack temporal learning capability essential for gait pattern analysis and fall risk prediction. Their accuracy (89.4\%/86.7\%) is substantially lower, and they cannot model the sequential progression of gait deterioration that precedes falls.

\textbf{CNN-GRU vs. Vision Transformer:} Vision Transformers achieve marginally higher accuracy (97.1\%) but require 67$\times$ more memory and 11.5$\times$ longer inference time, making them completely unsuitable for wearable embedded deployment. Their power consumption would drain a wearable battery within hours.

\subsection{Real-time Performance}

The system achieves the following real-time performance metrics on ESP32:
\begin{itemize}
    \item Sensor data acquisition: 10 ms per cycle
    \item Preprocessing (filtering + normalization): 5 ms
    \item CNN-GRU inference: 45 ms
    \item Total pipeline latency: 60 ms
    \item Alert generation: $<$100 ms from detection
    \item Cloud upload latency: 200--500 ms (WiFi dependent)
\end{itemize}

\subsection{Power Consumption Analysis}

\begin{itemize}
    \item ESP32 active mode: 160 mA
    \item MPU6050: 3.9 mA
    \item MAX30102: 0.7 mA (measurement mode)
    \item GPS (periodic): 45 mA (active), 0 mA (sleep)
    \item OLED display: 20 mA
    \item Total average: approximately 180 mA
    \item Battery life (1200 mAh): approximately 6.5 hours continuous, 18+ hours with duty cycling
\end{itemize}

\section{Advantages}

The Guardian 360 system offers the following advantages over existing solutions:

\begin{enumerate}
    \item \textbf{Predictive Capability:} Unlike reactive systems \cite{saraubon2018, kumar2024sensor}, Guardian 360 predicts fall risk before occurrence, enabling proactive intervention.
    \item \textbf{Lightweight Architecture:} The CNN-GRU model (180 KB) is deployable on low-cost ESP32 hardware without requiring cloud-based inference.
    \item \textbf{Multi-modal Monitoring:} Integrates motion analysis, physiological monitoring, and location tracking in a single wearable device.
    \item \textbf{Real-time Processing:} End-to-end pipeline latency of 60 ms enables immediate risk assessment and alert generation.
    \item \textbf{IoT Integration:} Firebase cloud connectivity enables remote monitoring, historical trend analysis, and multi-caregiver access.
    \item \textbf{User-Friendly Design:} Wrist-wearable form factor with OLED display and SOS button ensures elderly user acceptance.
    \item \textbf{Cost-Effective:} Total hardware cost under \$30 USD makes the system accessible for widespread deployment.
\end{enumerate}

\section{Limitations}

The current implementation has the following limitations:

\begin{enumerate}
    \item \textbf{Dataset Dependency:} The model is trained primarily on publicly available datasets with simulated falls; real-world elderly fall data remains limited.
    \item \textbf{Environmental Factors:} Extreme temperatures and moisture may affect sensor accuracy in outdoor environments.
    \item \textbf{WiFi Dependency:} Cloud features require WiFi connectivity; offline operation is limited to local alerts.
    \item \textbf{Battery Constraints:} Continuous high-frequency sensing limits battery life to 18--24 hours, requiring daily charging.
    \item \textbf{Individual Variability:} The model may require personalization for users with significantly atypical gait patterns due to specific medical conditions.
    \item \textbf{Limited Clinical Validation:} The system has not yet undergone large-scale clinical trials with elderly populations.
\end{enumerate}

\section{Future Scope}

The following directions are identified for future enhancement of Guardian 360:

\begin{itemize}
    \item \textbf{Tiny Transformers:} Exploration of lightweight attention mechanisms for improved temporal modeling within embedded constraints.
    \item \textbf{Vision-Augmented Posture Analysis:} Integration of low-power camera modules with pose estimation for comprehensive posture monitoring.
    \item \textbf{Multimodal AI Fusion:} Combining IMU data with audio (voice analysis for cognitive assessment) and environmental sensors for holistic health monitoring.
    \item \textbf{Cognitive Decline Detection:} Extending gait analysis to detect early signs of dementia and Alzheimer's disease, as gait deterioration correlates with cognitive decline \cite{tsiakiri2025}.
    \item \textbf{Emotion Recognition:} Integration of physiological signals for stress and emotional state monitoring \cite{vogka2025, rahmani2024}.
    \item \textbf{Federated Learning:} Privacy-preserving model improvement across multiple users without centralizing sensitive health data.
    \item \textbf{Smart Home Integration:} Connecting Guardian 360 with smart home systems for automated environmental adjustments (lighting, obstacle alerts).
    \item \textbf{Edge AI Optimization:} Neural architecture search (NAS) for discovering optimal CNN-GRU configurations for specific hardware targets.
    \item \textbf{5G/LoRa Connectivity:} Extending communication options for outdoor and rural deployment scenarios.
\end{itemize}

\section{Conclusion}

This paper presented Guardian 360, an AI-powered wearable system for elderly fall risk prediction using a novel lightweight CNN-GRU hybrid deep learning architecture. The system addresses the critical gap between reactive fall detection and proactive fall prevention by continuously analyzing gait patterns through wearable IMU sensors and predicting fall risk before occurrence.

The proposed CNN-GRU architecture achieves 96.2\% classification accuracy while maintaining an inference time of 45 ms on ESP32 hardware, demonstrating the feasibility of deploying sophisticated deep learning models on resource-constrained wearable devices. Compared to CNN-LSTM, the architecture reduces parameters by 47\% and memory requirements by 45\% while maintaining comparable accuracy, making it significantly more suitable for embedded deployment.

The integrated system combining gait analysis, fall prediction, physiological monitoring, GPS tracking, and IoT-enabled caregiver notification provides a comprehensive solution for elderly safety. By shifting the paradigm from reactive detection to predictive prevention, Guardian 360 enables timely intervention that can significantly reduce fall-related injuries and improve quality of life for elderly individuals.

Future work will focus on clinical validation with elderly populations, integration of lightweight attention mechanisms, and federated learning for privacy-preserving personalization, advancing toward truly intelligent and adaptive elderly healthcare systems.

\begin{thebibliography}{30}

\bibitem{who2024}
World Health Organization, ``Falls,'' WHO Fact Sheet, 2024. [Online]. Available: https://www.who.int/news-room/fact-sheets/detail/falls

\bibitem{salma2025}
I. Salma, ``Remote monitoring system for older adults at risk for complications: a scoping review,'' \textit{BMC Geriatrics}, vol. 25, 2025.

\bibitem{bandara2020}
A. Bandara, ``A sensor platform for non-invasive remote monitoring of older adults in real time,'' \textit{Academia}, 2020.

\bibitem{setu2025starlink}
S. S. Setu, ``A novel approach for a smart IoMT-based BAN for an old home healthcare monitoring system using Starlink,'' \textit{arXiv preprint}, arXiv:2512.22553, 2025.

\bibitem{setu2024}
S. S. Setu, ``Smart home solutions for elderly care,'' \textit{IEEE Xplore}, 2024.

\bibitem{augastiny2023}
R. Augastiny et al., ``IoT based health monitoring system for aged people using NodeMCU,'' \textit{IJNRD}, vol. 8, no. 9, 2023.

\bibitem{pasrija2025}
R. Pasrija et al., ``Enabling independence of elderly people using IoT technology,'' in \textit{Proc. IEEE Conf.}, 2025.

\bibitem{khalil2026}
R. A. Khalil et al., ``Redefining elderly care with agentic AI: challenges and opportunities,'' \textit{IEEE Access}, 2026.

\bibitem{ishaq2026}
M. Ishaq, ``A systematic review of fall detection and prediction technologies for older adults,'' \textit{Applied Sciences}, vol. 16, no. 4, p. 1929, 2026.

\bibitem{gaud2025}
N. Gaud, ``FIBiLS: fall detection of healthy elderly using IMU sensor and BiLSTM model,'' \textit{IEEE Sensors Journal}, 2025.

\bibitem{zhang2025}
Zhang et al., ``Artificial intelligence for elderly fall detection: state-of-the-art methods, applications, and challenges,'' \textit{ResearchGate}, 2025.

\bibitem{tsiakiri2025}
A. Tsiakiri, ``Wearable sensor technologies and gait analysis for early detection of dementia: trends and future directions,'' \textit{Sensors}, vol. 25, no. 24, p. 7669, 2025.

\bibitem{lim2024}
Z. K. Lim, ``Fall risk prediction using temporal gait features and machine learning approaches,'' \textit{Frontiers in Artificial Intelligence}, vol. 7, 2024.

\bibitem{li2024}
Li et al., ``Deep learning framework for fall detection,'' \textit{Cognitive Computation}, Springer, 2024.

\bibitem{kumar2024sensor}
Kumar et al., ``Sensor-based fall assistance system,'' \textit{Applied Sciences}, vol. 12, no. 7, p. 3276, 2024.

\bibitem{acm2024}
ACM Study, ``ML-based fall detection,'' \textit{PMC}, 2024.

\bibitem{attal2015}
F. Attal et al., ``Physical human activity recognition using wearable sensors,'' \textit{Sensors}, vol. 15, no. 12, pp. 29858--29894, 2015.

\bibitem{saraubon2018}
K. Saraubon, K. Anurugsa, and A. Kongsakpaibul, ``A smart system for elderly care using IoT and mobile technologies,'' in \textit{Proc. ACM Int. Conf.}, 2018.

\bibitem{alkhafajiy2019}
M. Al-khafajiy et al., ``Remote health monitoring of elderly through wearable sensors,'' \textit{Multimedia Tools and Applications}, Springer, 2019.

\bibitem{mdpi2025}
MDPI Sensors, ``Fall detection: systematic review,'' \textit{Sensors}, vol. 25, no. 21, p. 6540, 2025.

\bibitem{vogka2025}
N. Vogka, ``Smart wearable technologies for autonomous mental health monitoring in the elderly: a systematic review,'' \textit{EAI Endorsed Trans. Pervasive Health and Technology}, vol. 11, 2025.

\bibitem{rahmani2024}
M. H. Rahmani, ``EmoWear: wearable physiological and motion dataset for emotion recognition and context awareness,'' \textit{Scientific Data}, Nature, 2024.

\bibitem{deshmukh2026}
S. Deshmukh, ``An IoT-enabled smart healthcare monitoring system using machine learning for early health risk prediction,'' \textit{ResearchGate}, 2026.

\bibitem{harish2025}
Harish G, ``AI-powered Alzheimer's guider solution for patients with disease,'' 2025.

\bibitem{huang2024}
L. Huang et al., ``Automatic speech analysis for detecting cognitive decline of older adults,'' \textit{Frontiers in Public Health}, vol. 12, 2024.

\bibitem{sisfall}
A. Sucerquia, J. D. L\'{o}pez, and J. F. Vargas-Bonilla, ``SisFall: a fall and movement dataset,'' \textit{Sensors}, vol. 17, no. 1, p. 198, 2017.

\bibitem{mobiact}
G. Vavoulas et al., ``The MobiAct dataset: recognition of activities of daily living using smartphones,'' in \textit{Proc. ICT4AgeingWell}, 2016.

\bibitem{upfall}
L. Mart\'{i}nez-Villase\~{n}or et al., ``UP-Fall detection dataset: a multimodal approach,'' \textit{Sensors}, vol. 19, no. 9, p. 1988, 2019.

\bibitem{cho2014}
K. Cho et al., ``Learning phrase representations using RNN encoder-decoder for statistical machine translation,'' in \textit{Proc. EMNLP}, 2014.

\bibitem{tensorflow}
M. Abadi et al., ``TensorFlow: a system for large-scale machine learning,'' in \textit{Proc. OSDI}, 2016.

\end{thebibliography}

\end{document}

